% !TeX program = xelatex
% Bao cao theo ma nguon hien tai va 6 log benchmark ngay 17/10/2026.
% Uu tien XeLaTeX; cung ho tro LuaLaTeX va pdfLaTeX co VnTeX.
% Moi so do va bieu do nam trong file nay.
\documentclass[11pt,a4paper]{article}
\usepackage{iftex}
\ifPDFTeX
  % pdfLaTeX khong ho tro fontspec; dung ma hoa tieng Viet cua VnTeX.
  \usepackage[utf8]{inputenc}
  \usepackage[T5]{fontenc}
  \usepackage[vietnamese]{babel}
\else
  \usepackage{fontspec}
  \defaultfontfeatures{Ligatures=TeX}
  \IfFontExistsTF{Times New Roman}
    {\setmainfont{Times New Roman}}
    {\IfFontExistsTF{TeX Gyre Termes}
      {\setmainfont{TeX Gyre Termes}}
      {\setmainfont{Latin Modern Roman}}}
  \IfFontExistsTF{TeX Gyre Heros}
    {\setsansfont{TeX Gyre Heros}}
    {\IfFontExistsTF{Arial}
      {\setsansfont{Arial}}
      {\setsansfont{Latin Modern Sans}}}
  \IfFontExistsTF{Latin Modern Mono}
    {\setmonofont{Latin Modern Mono}[Scale=MatchLowercase]}
    {\setmonofont{Consolas}[Scale=MatchLowercase]}
\fi
\usepackage[top=1.7cm,bottom=1.8cm,left=2cm,right=2cm]{geometry}
\usepackage{amsmath,booktabs,tabularx,array,enumitem}
\usepackage{xcolor,tikz,pgfplots,caption,fancyhdr,xurl}
\usepackage[unicode,hidelinks]{hyperref}
\urlstyle{same}
\usetikzlibrary{arrows.meta,positioning,calc}
\usepgfplotslibrary{groupplots}
\pgfplotsset{compat=1.18}
\definecolor{accent}{HTML}{175B78}
\definecolor{second}{HTML}{CE7B35}
\definecolor{ink}{HTML}{18333F}
\definecolor{muted}{HTML}{53636B}
\definecolor{rulegray}{HTML}{D5DEE3}
\hypersetup{
  pdftitle={Báo cáo hệ thống phân tích video đa camera},
  pdfauthor={},
  pdfsubject={Kiến trúc, vận hành và benchmark shared / per_camera}
}
\renewcommand{\figurename}{Hình}
\renewcommand{\tablename}{Bảng}
\renewcommand{\arraystretch}{1.18}
\setlength{\parindent}{0pt}
\setlength{\parskip}{5pt}
\setlength{\emergencystretch}{2em}
\setlist[itemize]{leftmargin=1.4em,itemsep=3pt,topsep=4pt}
\setlist[enumerate]{leftmargin=1.6em,itemsep=4pt,topsep=4pt}
\captionsetup{font=small,labelfont=bf,skip=6pt}
\newcolumntype{Y}{>{\raggedright\arraybackslash}X}
\newcommand{\code}[1]{{\rmfamily\detokenize{#1}}}
\newcommand{\shared}{\code{shared}}
\newcommand{\percam}{\code{per_camera}}
\newcommand{\source}[1]{\par{\footnotesize\color{muted}\textit{Mã nguồn:} #1}\par}
\newcommand{\insight}[1]{%
  \par\smallskip\begingroup\setlength{\fboxsep}{8pt}%
  \noindent\fcolorbox{accent!25}{accent!5}{%
    \parbox{\dimexpr\linewidth-2\fboxsep-2\fboxrule\relax}{#1}}%
  \endgroup\par\smallskip}
\pagestyle{fancy}
\fancyhf{}
\fancyhead[L]{\small\color{muted}HỆ THỐNG PHÂN TÍCH VIDEO ĐA CAMERA}
\fancyhead[R]{\small\color{muted}Báo cáo kỹ thuật}
\fancyfoot[L]{\footnotesize\color{muted}Mã nguồn và log trong dự án}
\fancyfoot[R]{\small\thepage}
\renewcommand{\headrulewidth}{0.3pt}
\setlength{\headheight}{15pt}
\tikzset{
  block/.style={draw=accent!65,fill=accent!4,rounded corners=3pt,
    align=center,minimum height=0.72cm,inner sep=5pt,font=\small},
  flow/.style={-{Stealth[length=2mm]},thick,draw=accent},
  side/.style={draw=rulegray,fill=white,rounded corners=3pt,
    align=left,inner sep=7pt,font=\footnotesize}
}
\pgfplotsset{
  reportbar/.style={
    ybar,bar width=17pt,
    width=0.455\textwidth,height=6.55cm,
    symbolic x coords={1,4,8},xtick=data,
    xlabel={Số camera},ylabel near ticks,
    ymin=0,enlarge x limits=0.23,
    axis line style={draw=rulegray},
    tick style={draw=rulegray},
    ymajorgrids,grid style={draw=rulegray!70},
    tick label style={font=\small},
    label style={font=\small},
    title style={font=\small\bfseries,align=center},
    nodes near coords,
    nodes near coords style={font=\scriptsize,anchor=south},
    every node near coord/.append style={/pgf/number format/fixed,
      /pgf/number format/precision=1},
    legend style={draw=none,font=\small,legend columns=2,
      at={(0.5,1.13)},anchor=south},
    legend cell align=left
  }
}

\begin{document}
\thispagestyle{plain}
\begin{center}
  {\Large\bfseries\color{ink}
    BÁO CÁO HỆ THỐNG PHÂN TÍCH VIDEO ĐA CAMERA}\\[5pt]
  {\normalsize Kiến trúc, phương thức vận hành và so sánh hiệu năng}\\[3pt]
  {\small Ngày lập: 18/10/2026 \quad | \quad Log benchmark: 17/10/2026}
\end{center}

\section{Tổng quan hệ thống}
Hệ thống nhận hình ảnh từ nhiều camera, phát hiện đối tượng bằng YOLO,
gán ID theo dõi bằng ByteTrack và phát hiện hai loại sự kiện:
\textbf{vượt vạch} và \textbf{đi vào vùng cấm}.
Sự kiện được lưu vào SQLite cùng ảnh minh chứng; video có chú thích được
hiển thị tại máy và/hoặc phát qua RTSP. API phục vụ theo dõi trạng thái,
bật/tắt camera và tra cứu lịch sử.

\begin{center}
\begin{tikzpicture}[
  node distance=4.5mm,
  overview/.style={block,text width=5.6cm,
    font=\fontsize{9.5pt}{11.5pt}\selectfont,
    inner xsep=6pt,inner ysep=5pt,
    execute at begin node={\setlength{\parskip}{0pt}}}
]
% Dat theo bien hop de khoang cach khong phu thuoc chieu cao chu/font.
\node[overview] (manager)
  {\textbf{SystemManager}\\Đọc cấu hình, khởi tạo và điều phối};
\node[overview,below=of manager] (capture)
  {\textbf{CameraWorker $\times N$}\\Nhận frame mới nhất của từng camera};
\node[overview,below=of capture] (infer)
  {\textbf{InferenceWorker}\\YOLO: bbox, class, confidence};
\node[overview,below=of infer] (track)
  {\textbf{Tracker}\\ByteTrack riêng cho từng camera};
\node[overview,below=of track] (event)
  {\textbf{EventEngine}\\Vượt vạch / đi vào vùng cấm};
\draw[flow] (manager.south)--(capture.north);
\draw[flow] (capture.south)--(infer.north);
\draw[flow] (infer.south)--(track.north);
\draw[flow] (track.south)--(event.north);
\node[overview,text width=5.4cm,below=10mm of event.south,xshift=-3.4cm] (writer)
  {\textbf{WriterWorker $\rightarrow$ EventWriter}\\Ghi sự kiện bất đồng bộ};
\node[overview,text width=5.4cm,below=10mm of event.south,xshift=3.4cm] (drawing)
  {\textbf{Drawing Utils}\\Vẽ bbox, ID, vùng, vạch và FPS};
\coordinate (fork) at ($(event.south)+(0,-5mm)$);
\draw[thick,draw=accent] (event.south)--(fork);
\draw[flow] (fork)-|(writer.north);
\draw[flow] (fork)-|(drawing.north);
\node[overview,text width=5.4cm,below=5mm of writer] (db)
  {\textbf{SQLite + ảnh minh chứng}\\Ảnh JPEG Base64 trong bảng events};
\node[overview,text width=5.4cm,below=5mm of drawing] (stream)
  {\textbf{StreamWorker $\rightarrow$ FFmpeg}\\H.264 $\rightarrow$ RTSP server};
\draw[flow] (writer.south)--(db.north);
\draw[flow] (drawing.south)--(stream.north);
\node[side,text width=3.5cm,inner sep=6pt,right=13mm of manager] (api)
  {\textbf{FastAPI}\\Trạng thái, bật/tắt camera\\và truy vấn SQLite};
\draw[thick,draw=accent,{Stealth[length=2mm]}-{Stealth[length=2mm]}]
  (manager.east)--(api.west);
\end{tikzpicture}
\captionof{figure}{Workflow tổng thể. Mũi tên chính biểu diễn thứ tự xử lý;
Manager gọi các bước inference, tracking và sự kiện.}
\end{center}

\begin{tabularx}{\linewidth}{@{}p{3.3cm}Y@{}}
\toprule
\textbf{Đầu vào / đầu ra} & \textbf{Ý nghĩa}\\
\midrule
Đầu vào & RTSP/RTSPS, webcam hoặc video file; mỗi camera có cấu hình riêng.\\
Nhánh dữ liệu & Lưu thời gian, camera, track ID, loại sự kiện và ảnh gốc.\\
Nhánh video & Xem video đã vẽ qua preview hoặc RTSP; server RTSP chạy bên ngoài.\\
Cách tổ chức & Capture theo từng camera; xử lý trung tâm; ghi DB và phát stream ở luồng riêng.\\
\bottomrule
\end{tabularx}

\smallskip
Hệ thống theo dõi ID \textbf{trong từng camera}; chưa hợp nhất danh tính giữa
các camera. Camera không cấu hình vạch hoặc vùng vẫn nhận diện, theo dõi
và xuất video.
\source{\path{src/manager.py}; \path{src/schemas.py}; \path{src/api.py}.}

\newpage
\section{Phương thức vận hành}
\subsection{Luồng CameraWorker}
Mỗi camera có một \code{CameraWorker} và một luồng thu nhận riêng.
Nhờ đó việc chờ đọc hình của camera này không trực tiếp chặn thao tác
đọc của các camera còn lại.

\begin{enumerate}
\item \textbf{Mở nguồn.} OpenCV mở webcam, video file hoặc RTSP. Với RTSP,
timeout mở và đọc được đặt là 3 giây. Video file được đọc theo FPS nguồn
để tránh chạy nhanh hơn thời gian thực.
\item \textbf{Đóng gói frame.} Mỗi \code{FramePacket} gồm
\code{camera_id}, ảnh gốc, thời điểm nhận \code{timestamp} và
\code{generation} nhận diện lần mở nguồn.
\item \textbf{Giữ frame mới nhất.} Hàng đợi chỉ có một vị trí. Khi frame mới
đến mà hàng đợi đầy, worker thay frame cũ. Manager lấy frame hiện có,
không chờ camera chậm và không chờ đủ batch.
\item \textbf{Kết nối lại.} Khi đọc thất bại, worker xóa frame chờ, đóng nguồn
và đợi 2 giây trước khi thử lại. Manager bỏ frame quá cũ; khi nguồn được mở
lại hoặc gián đoạn vượt ngưỡng mặc định 3 giây, tracker và trạng thái ID
của camera được reset.
\end{enumerate}

\insight{\textbf{Mục đích của hàng đợi một frame:} ưu tiên hình ảnh mới và
hạn chế tích lũy độ trễ khi xử lý chậm. Đổi lại, hệ thống có thể bỏ qua
chuyển động hoặc sự kiện xảy ra giữa những frame không được xử lý.}
\source{\path{src/camera_worker.py}; \path{src/manager.py}, hàm
\code{_collect_frames()}.}

\subsection{Luồng InferenceWorker $\rightarrow$ Tracker}
Manager thu tối đa một frame mới từ mỗi camera trong một lượt, tổng số
không vượt \code{batch_size} (cấu hình hiện tại: 8).
Khi số camera lớn hơn giới hạn này, điểm bắt đầu thu được luân phiên để
các camera đều có cơ hội được xử lý.

\begin{tabularx}{\linewidth}{@{}p{3.1cm}Y@{}}
\toprule
\textbf{Bước} & \textbf{Dữ liệu và xử lý}\\
\midrule
InferenceWorker & YOLO nhận danh sách ảnh, áp dụng weights, confidence,
classes, kích thước ảnh và device trong cấu hình; trả bbox dạng
$(x_1,y_1,x_2,y_2)$, confidence và class ID.\\
DetectionPacket & Giữ nguyên camera, timestamp, generation và frame gốc;
bổ sung kết quả phát hiện.\\
Tracker & Chuyển detection sang ByteTrack của đúng camera, lọc track chưa
hợp lệ; trả \code{TrackingPacket} có thêm \code{track_ids}.\\
\bottomrule
\end{tabularx}

\smallskip
Hai cách bố trí model dùng chung phần capture, tracking, xử lý sự kiện
và đầu ra:
\begin{itemize}
\item \textbf{\shared:} một model YOLO nhận batch gồm frame của nhiều camera.
\item \textbf{\percam:} mỗi camera có một model YOLO riêng, nhưng Manager
gọi các model \textbf{lần lượt}. Vì mỗi lượt chỉ lấy tối đa một frame/camera,
batch thực tế của từng model là 1.
\end{itemize}

Inference, tracking và EventEngine được gọi trong cùng vòng xử lý Manager.
Toàn bộ inference của lượt hiện tại hoàn tất trước khi Manager chạy
tracking và sự kiện cho các kết quả. Vì vậy, ở \percam, một camera còn
chịu thời gian chờ các lần inference khác trong cùng lượt.

ByteTrack được quản lý theo \code{camera_id}; cùng một số ID ở hai camera
không có nghĩa là cùng một đối tượng ngoài thực tế.
\source{\path{src/inference_worker.py}; \path{src/tracker.py};
\path{src/manager.py}, hàm \code{_predict()} và \code{run()}.}

\newpage
\subsection{Logic xử lý sự kiện}
Mỗi camera có trạng thái sự kiện độc lập: phía vạch trước đó, tâm trước đó,
trạng thái trong vùng và thời điểm cuối cùng nhìn thấy từng track ID.
Điểm đại diện đối tượng là \textbf{tâm bounding box}:
\[
 C=\left(\frac{x_1+x_2}{2},\frac{y_1+y_2}{2}\right).
\]
Vạch và đa giác được kiểm tra trên frame gốc. Khi chọn
\code{coordinates: normalized}, tọa độ trong khoảng $[0,1]$ được nhân
với chiều rộng và chiều cao của frame trước khi tính.

\textbf{Phát hiện vượt vạch.}
Với vạch có thứ tự $A\rightarrow B$, phía của tâm $P$ được xác định bởi:
\[
 s(P)=(B_x-A_x)(P_y-A_y)-(B_y-A_y)(P_x-A_x).
\]
Một track được tính vượt vạch khi tâm cũ và tâm mới đổi phía
($s(C_{\text{cũ}})\,s(C_{\text{mới}})<0$), đồng thời đoạn nối hai tâm
giao với \textbf{đoạn vạch hữu hạn} $AB$. Điều kiện thứ hai loại trường hợp
chỉ đi qua phần kéo dài ngoài hai đầu vạch.

\begin{center}
\begin{tikzpicture}[scale=0.9]
\draw[line width=1.3pt,accent,-{Stealth}] (-2.5,0)--(2.5,0);
\node[below] at (-2.5,0) {$A$};
\node[below] at (2.5,0) {$B$};
\fill[second] (-0.9,-1.05) circle (2.7pt);
\fill[accent] (0.85,1.05) circle (2.7pt);
\draw[flow,draw=second] (-0.9,-1.05)--(0.85,1.05);
\node[left,font=\small] at (-0.9,-1.05) {$C_{\text{cũ}}$, $s<0$};
\node[right,font=\small] at (0.85,1.05) {$C_{\text{mới}}$, $s>0$};
\node[font=\small,align=center] at (6,0)
  {$s<0 \rightarrow s>0$: \code{left_to_right}\\
   $s>0 \rightarrow s<0$: \code{right_to_left}};
\end{tikzpicture}
\captionof{figure}{Quy ước hướng theo dấu tích có hướng.
Đảo thứ tự $A,B$ sẽ đảo tên hướng.}
\end{center}

Hướng dương tăng \code{right_count}; hướng âm tăng \code{left_count}.
Tên trái/phải là quy ước theo $A,B$, không cố định theo trái/phải màn hình.
Nếu tâm nằm đúng trên vạch, code giữ phía và tâm khác 0 gần nhất để
so sánh ở lần quan sát tiếp theo.

\textbf{Phát hiện vào vùng cấm.}
Hệ thống dùng \code{pointPolygonTest} để kiểm tra tâm; cả điểm trên biên
cũng được xem là trong vùng.
\[
 \text{zone\_entry}=\text{inside}_{\text{hiện tại}}
 \land \neg\,\text{inside}_{\text{trước}}.
\]

\begin{tabularx}{\linewidth}{@{}p{6.6cm}Y@{}}
\toprule
\textbf{Tình huống quan sát} & \textbf{Kết quả}\\
\midrule
Lần đầu thấy ID và tâm đã ở trong vùng & Sinh một \code{zone_entry}.\\
ID tiếp tục ở trong vùng qua nhiều frame & Giữ cảnh báo, không ghi lặp entry.\\
ID ra ngoài, sau đó đi vào lại & Sinh entry mới.\\
ID vượt vạch và vào vùng trong cùng frame & Có thể sinh cả hai loại sự kiện.\\
\bottomrule
\end{tabularx}

\smallskip
Trạng thái ID không xuất hiện quá 60 giây được dọn. Reset sau gián đoạn
xóa trạng thái ID, giữ tổng đếm vạch trong phiên; chưa khôi phục tổng đếm
sau khi khởi động lại chương trình. Vạch và vùng được bật độc lập theo
cấu hình từng camera.

\insight{\textbf{Giới hạn logic:} kết quả phụ thuộc vào ID và các tâm đã
quan sát. Đổi ID, rung vị trí sát biên hoặc bỏ frame có thể ảnh hưởng số
sự kiện; code chưa có cơ chế chống rung hình học riêng.}
\source{\path{src/event_engine.py}; \path{src/schemas.py};
\path{src/config.py}.}

\newpage
\subsection{Logic EventWriter và API truy xuất}
Manager chỉ gửi kết quả có \code{zone_entry} hoặc vượt vạch vào
\code{WriterWorker}. Luồng writer xử lý hàng đợi hữu hạn, mặc định 128
\textbf{frame có sự kiện}, rồi gọi \code{EventWriter.save()}.

\textbf{Quy trình lưu:}
lọc các sự kiện cần ghi $\rightarrow$ mã hóa ảnh gốc thành JPEG chất lượng 80
$\rightarrow$ chuyển Base64 $\rightarrow$ chèn SQLite.
Bảng \code{events} gồm \code{id}, \code{timestamp}, \code{camera_id},
\code{track_id}, \code{event_type} và \code{snapshot_base64}.
Một frame có thể tạo nhiều bản ghi; các bản ghi đó được commit trong cùng
một transaction và dùng cùng nội dung ảnh.

Ảnh minh chứng là \textbf{frame gốc chưa vẽ}, được lưu ngay trong database.
Thời gian sự kiện lấy bằng \code{datetime.now()} khi xử lý.
SQLite dùng WAL và kết nối ghi thuộc riêng luồng writer.
Khi dừng bình thường, writer ghi hết phần hàng đợi đã nhận.
Nếu hàng đợi đầy hoặc writer lỗi, pipeline báo lỗi và dừng; sự kiện
không enqueue được khi đầy chưa được lưu.

\begin{tabularx}{\linewidth}{@{}p{6.2cm}Y@{}}
\toprule
\textbf{API} & \textbf{Vai trò}\\
\midrule
\code{GET /status} & Trạng thái pipeline, model mode, FPS và hàng đợi.\\
\code{GET /cameras} & Danh sách và trạng thái các camera.\\
\code{GET /cameras/{camera_id}} & Trạng thái và số liệu một camera.\\
\code{PATCH /cameras/{camera_id}} & Bật/tắt camera bằng \code{enabled}.\\
\code{GET /events} & Lọc lịch sử theo camera, thời gian; phân trang.\\
\bottomrule
\end{tabularx}

\smallskip
Mỗi truy vấn \code{/events} mở một \code{EventWriter} với kết nối riêng,
đọc rồi đóng ngay trong cùng thread. API trả \code{items}, \code{total},
\code{limit}, \code{offset}; mỗi item chứa metadata, không chứa ảnh Base64.
Kết quả sắp xếp mới nhất trước; \code{limit} mặc định 100, tối đa 500.
Thời gian lọc dùng giờ máy không kèm múi giờ và bao gồm hai đầu mút.
Ví dụ: \path{/events?camera_id=cam_01&limit=20&offset=0}.

API sử dụng cùng Manager; lệnh bật/tắt được đưa về hàng đợi điều khiển
để Manager thực hiện. Phương thức nội bộ \code{EventWriter.get(id)} có thể
đọc ảnh, nhưng API hiện chưa cung cấp endpoint lấy riêng ảnh sự kiện.
\source{\path{src/output_workers.py}; \path{src/event_writer.py};
\path{src/api.py}.}

\subsection{Luồng stream output}
\code{draw_results()} vẽ trên bản sao của frame: bbox, track ID, class ID,
confidence, vạch, đa giác, tổng đếm và FPS. Bbox màu đỏ khi tâm nằm trong
vùng cấm, màu xanh khi ở ngoài. Cùng kết quả này có thể được dùng cho
preview và nhánh RTSP.

Mỗi camera có đầu ra RTSP dùng một \code{StreamWorker}, hàng đợi một frame
và một tiến trình FFmpeg. Worker ưu tiên frame mới, giới hạn tốc độ gửi;
frame bị thay thế được đếm riêng. Trước khi gửi, ảnh được resize giữ tỷ lệ
và thêm nền đen theo kích thước output (mặc định từ YAML:
$800\times450$, tối đa 15 FPS).

FFmpeg nhận BGR24 qua stdin, mã hóa H.264 bằng \code{libx264}
(preset \code{ultrafast}, tune \code{zerolatency}), rồi publish RTSP/TCP
đến server ngoài chương trình. Pipe lỗi sẽ được mở lại; watchdog kết thúc
FFmpeg nếu thao tác ghi bị treo 3 giây.

\insight{\textbf{Cách đọc số liệu đầu ra:}
15 FPS là giới hạn gửi, không bảo đảm luôn đạt được.
Chỉ số \code{e2e_ms} kết thúc khi Python ghi/flush xong sang FFmpeg;
chưa đo hoàn tất encode, mạng, buffer RTSP hoặc màn hình người xem.}
\source{\path{src/utils.py}; \path{src/output_workers.py};
\path{src/stream_ouput.py}.}

\newpage
\section{Kết quả benchmark}
\subsection*{Mục tiêu và hai phương án so sánh}
Benchmark so sánh \textbf{hai cách bố trí model trong pipeline hiện có}:
\shared{} dùng một model cho batch nhiều camera;
\percam{} dùng model riêng từng camera và gọi tuần tự.
Các kịch bản gồm 1, 4 và 8 camera, mỗi chế độ có một lượt chạy thành công.
Mục tiêu là đối chiếu thông lượng xử lý và độ trễ khi tăng số camera.

Dữ liệu lấy từ sáu log ngày 17/10/2026 tại
\path{benchmark/local/1cam/}, \path{benchmark/local/4cam/},
\path{benchmark/local/8cam/}; mỗi thư mục có \code{share.log} và
\code{per.log}. Log ghi nhận chạy local trên Windows, \code{cuda:0},
\code{batch_size=8}, bật preview và streaming; FFmpeg xuất
$800\times450$ với giới hạn 15 FPS. Log chưa lưu đầy đủ weights, phiên bản
thư viện, cấu hình model và mức sử dụng CPU/GPU/RAM/VRAM của mỗi phiên.

\subsection*{Cách chọn khoảng đo và tính chỉ số}
Lấy snapshot định kỳ thứ 3 làm mốc đầu và thứ 9 làm mốc cuối của từng
phiên thành công, tương ứng khoảng giây 90--272 sau khởi động và khoảng
180 giây đo/cấu hình. Phần đầu khởi động được loại khỏi phép tổng hợp;
tiền tố lượt chạy lỗi trong \path{8cam/per.log} cũng được loại.

\[
 F_{\text{tổng}}=\frac{
     P_{\text{cuối}}-P_{\text{đầu}}}{
     t_{\text{cuối}}-t_{\text{đầu}}},\qquad
 \overline F_{\mathrm{camera}}=\frac{F_{\text{tổng}}}{N}.
\]
$P$ là bộ đếm \code{processed_frames}; ví dụ 4 camera \shared{}:
$9\,374/180{,}146=52{,}04$ FPS.
FPS xuất tính tương tự từ chênh lệch \code{sent_frames}, rồi chia số camera.

\begin{itemize}
\item $L_{\text{xử lý}}$: từ lúc Python nhận frame đến khi hoàn tất
tracking, xử lý sự kiện và gửi vào hàng đợi writer; trước bước vẽ.
\item $L_{\mathrm{pipe}}$: từ lúc nhận frame đến khi Python ghi/flush
xong sang FFmpeg, tức trường \code{e2e_ms}.
\end{itemize}
Hai cột độ trễ dưới đây là trung bình các \code{avg_ms} ở snapshot thứ
4--9 của tất cả camera, mỗi camera và mỗi snapshot có trọng số như nhau.
Mỗi snapshot phản ánh cửa sổ gần nhất 5 giây; đây \textbf{không phải}
trung bình toàn bộ frame trong khoảng 180 giây hoặc độ trễ phía người xem.

\begin{center}
\small
\setlength{\tabcolsep}{5pt}
\begin{tabular}{@{}clrrrrr@{}}
\toprule
\shortstack{Số\\camera} & Chế độ &
\shortstack{FPS\\tổng} &
\shortstack{FPS/camera\\trung bình} &
\shortstack{$L_{\text{xử lý}}$\\(ms)} &
\shortstack{$L_{\mathrm{pipe}}$\\(ms)} &
\shortstack{FPS xuất\\TB/camera}\\
\midrule
1 & \shared & 28,61 & 28,61 & 37,46 & 61,28 & 12,71\\
1 & \percam & 28,58 & 28,58 & 37,71 & 63,82 & 12,72\\
\addlinespace
4 & \shared & \textbf{52,04} & \textbf{13,01} & \textbf{71,33} & 113,74 & 12,77\\
4 & \percam & 36,16 & 9,04 & 102,58 & \textbf{107,30} & 9,04\\
\addlinespace
8 & \shared & \textbf{29,26} & \textbf{3,66} & \textbf{212,64} & \textbf{220,84} & 3,66\\
8 & \percam & 24,61 & 3,08 & 286,88 & 293,27 & 3,08\\
\bottomrule
\end{tabular}
\captionof{table}{Tổng hợp từ sáu log. FPS cao hơn tốt hơn; độ trễ thấp hơn tốt hơn.}
\end{center}

\insight{\textbf{Lưu ý khi so sánh:} \code{inference_ms} của \shared{}
là thời gian cho cả batch, còn của \percam{} là thời gian một model/camera.
Báo cáo ưu tiên FPS và độ trễ xử lý để tránh so trực tiếp hai lượng công
việc khác nhau.}

\newpage
\subsection*{Biểu đồ so sánh}
Màu xanh biểu diễn \shared{}, màu cam biểu diễn \percam{}.
Nhãn cột được làm tròn một chữ số thập phân; bảng số liệu giữ hai chữ số.

\begin{center}
\begin{tikzpicture}
\begin{groupplot}[
  reportbar,
  group style={group size=2 by 2,horizontal sep=1.35cm,vertical sep=2.3cm}
]
\nextgroupplot[title={(a) Thông lượng toàn hệ thống},
  ylabel={FPS tổng},ymax=65,ytick={0,10,20,30,40,50,60}]
\addplot[fill=accent,draw=accent] coordinates {(1,28.60539) (4,52.03557) (8,29.26313)};
\addplot[fill=second,draw=second] coordinates {(1,28.58063) (4,36.15781) (8,24.61433)};
\legend{shared,per\_camera}

\nextgroupplot[title={(b) Tốc độ xử lý mỗi camera},
  ylabel={FPS/camera trung bình},ymax=36,ytick={0,5,10,15,20,25,30}]
\addplot[fill=accent,draw=accent] coordinates {(1,28.60539) (4,13.00889) (8,3.65789)};
\addplot[fill=second,draw=second] coordinates {(1,28.58063) (4,9.03945) (8,3.07679)};
\draw[dashed,muted,thick] (rel axis cs:0,0.416667)--(rel axis cs:1,0.416667)
  node[pos=0.97,above,font=\scriptsize,fill=white,inner sep=1pt] {15 FPS};

\nextgroupplot[title={(c) Độ trễ xử lý},
  ylabel={ms},ymax=350,ytick={0,50,100,150,200,250,300}]
\addplot[fill=accent,draw=accent] coordinates {(1,37.455) (4,71.33417) (8,212.63583)};
\addplot[fill=second,draw=second] coordinates {(1,37.71333) (4,102.57667) (8,286.87563)};

\nextgroupplot[title={(d) Độ trễ nhận frame $\rightarrow$ pipe FFmpeg},
  ylabel={ms},ymax=350,ytick={0,50,100,150,200,250,300}]
\addplot[fill=accent,draw=accent] coordinates {(1,61.28333) (4,113.73917) (8,220.83875)};
\addplot[fill=second,draw=second] coordinates {(1,63.815) (4,107.30083) (8,293.27229)};
\end{groupplot}
\end{tikzpicture}
\captionof{figure}{So sánh hai chế độ tại 1, 4 và 8 camera.
Đường 15 FPS là mốc tham chiếu, không phải tốc độ inference cố định hay
yêu cầu đã được xác nhận. Không có thanh sai số vì chỉ có một lượt chạy/cấu hình.}
\end{center}

\insight{\textbf{Điểm nổi bật:}
ở 4 camera, \shared{} xử lý được 52,04 FPS tổng so với 36,16 FPS,
tăng 43,9\%. Ở 8 camera, \shared{} vẫn tốt hơn \percam{}, nhưng chỉ đạt
trung bình 3,66 FPS/camera.}

\textbf{Cách đọc biểu đồ:}
(a) cho biết năng lực xử lý tổng;
(b) cho biết phần năng lực trung bình mỗi camera nhận được;
(c) phản ánh tốc độ hoàn thành phân tích;
(d) bổ sung thời gian đến nhánh xuất video.
Không dùng (d) để suy ra độ trễ phát hình RTSP ở máy người xem.

\newpage
\section{Nhận xét, đánh giá và lựa chọn}
\textbf{Lựa chọn đề xuất: dùng \shared{} làm mặc định cho nhiều camera
khi cùng sử dụng một bộ weights và cấu hình model.}
Trong các lượt đo hiện có, phương án này tăng thông lượng và giảm độ trễ
phân tích ở cả 4 và 8 camera. Cấu trúc cũng chỉ cần nạp một model;
mức tiết kiệm RAM/VRAM cụ thể chưa được đo.

\begin{tabularx}{\linewidth}{@{}p{2.3cm}Y@{}}
\toprule
\textbf{Kịch bản} & \textbf{Nhận xét từ số liệu}\\
\midrule
1 camera & 28,61 so với 28,58 FPS: chênh lệch khoảng 0,09\%.
Hai chế độ gần tương đương trong lượt đo này; chưa đủ để kết luận
có ưu thế đáng kể.\\
4 camera & \shared{} tăng FPS tổng \textbf{43,9\%}, giảm độ trễ xử lý
\textbf{30,5\%}. Tốc độ trung bình đạt 13,01 FPS/camera,
so với 9,04 FPS/camera.\\
8 camera & \shared{} tăng FPS tổng \textbf{18,9\%}, giảm độ trễ xử lý
\textbf{25,9\%} và giảm độ trễ đến pipe \textbf{24,7\%}.
Tuy vậy, 3,66 FPS/camera còn thấp nếu cần hình ảnh phân tích liên tục.\\
\bottomrule
\end{tabularx}

\subsection*{Vì sao không chỉ nhìn FPS tổng?}
\begin{itemize}
\item \textbf{Nhánh xuất có thể làm tăng trễ.}
Với 4 camera, dù \shared{} xử lý nhanh hơn, độ trễ đến pipe tăng
từ 107,30 lên 113,74 ms, khoảng 6,0\%.
Các snapshot ghi nhận thời gian chờ hàng đợi đầu ra trung bình
36,89 ms so với 0,33 ms; có 169 frame xuất bị thay thế trong khoảng đo
của \shared{}, so với 0 của \percam{}.
\item \textbf{FPS trung bình chưa bảo đảm mọi camera đều đạt.}
Ở 8 camera \shared{}, từng camera đạt khoảng 3,40--3,86 FPS.
Nếu đặt mục tiêu 15 FPS xử lý cho mỗi camera thì cả hai chế độ ở
4 và 8 camera đều chưa đạt trong các lượt đo này.
\item \textbf{Bỏ frame đầu ra khác bỏ sót sự kiện.}
Ở 1 camera, xử lý khoảng 28,6 FPS nhưng đầu ra chỉ khoảng 12,7 FPS;
nhiều frame bị thay ở nhánh stream. Số này không đo tỷ lệ đếm sai
hoặc sự kiện bị bỏ sót.
\end{itemize}

Trong các khoảng được chọn, bộ đếm reconnect và lỗi gửi không tăng;
hàng đợi sự kiện bằng 0 tại các snapshot. Điều đó mô tả trạng thái tại
các mốc đo, chưa chứng minh hàng đợi luôn rỗng hoặc hệ thống không thể
mất sự kiện.

\subsection*{Phạm vi kết luận và bước kiểm chứng tiếp theo}
Mỗi cấu hình mới có một lượt đo, chưa có khoảng tin cậy và nhãn đối chứng
độ chính xác. FPS nguồn giữa các camera khác nhau; log cũng chưa lưu đầy đủ
tham số model và phiên bản. Vì vậy, các chênh lệch trên mô tả
\textbf{những phiên đã đo}, chưa tách được ảnh hưởng của mọi yếu tố.
So sánh ưu tiên hai chế độ trong cùng số camera; không suy ra
khả năng tăng camera theo tỷ lệ tuyến tính hay kết luận Docker nhanh hơn.

Để chốt cấu hình triển khai, cần lặp ít nhất 3 lượt/chế độ với cùng video,
weights, tham số, preview và đầu ra; đổi thứ tự chạy; đo thêm CPU/GPU,
RAM/VRAM, FPS camera chậm nhất và tỷ lệ phát hiện sự kiện đúng.
Hiện chưa đủ bằng chứng để quy nút thắt ở 8 camera cho riêng GPU,
decode, tracking hay encode.

\insight{\textbf{Quyết định từ dữ liệu hiện có:}
chọn \shared{} cho pipeline nhiều camera; giữ \percam{} làm phương án
đối chứng. Số camera triển khai phải được quyết định theo FPS và độ trễ
ứng dụng cần, rồi đo lại trên nguồn và phần cứng thực tế.}

{\footnotesize\color{muted}
\textbf{Nguồn đối chiếu:}
mã nguồn trong \path{src/}, cấu hình trong \path{configs/},
\path{docs/benchmark-models.md}, sáu log \path{benchmark/local/}.
Chi tiết khoảng đo được ghi trong \path{benchmark-evidence.csv}
đi kèm mã LaTeX. Báo cáo không dùng kết quả smoke test để thay thế benchmark.}

% Provenance: snapshot dinh ky #3 -> #9 cua phien thanh cong.
% Latency: cac snapshot #4 -> #9, avg_ms cua moi camera dong trong so.
% 1cam/shared 16:27:48.930 -> 16:30:48.966; 5150 frame / 180.036 s.
% 1cam/per    16:21:35.354 -> 16:24:35.406; 5146 frame / 180.052 s.
% 4cam/shared 16:46:11.942 -> 16:49:12.088; 9374 frame / 180.146 s.
% 4cam/per    16:40:44.899 -> 16:43:45.192; 6519 frame / 180.293 s.
% 8cam/shared 17:14:23.936 -> 17:17:24.402; 5281 frame / 180.466 s.
% 8cam/per    17:05:31.300 -> 17:08:32.089; 4450 frame / 180.789 s.
\end{document}
